% Provenance (never printed): see sources/outline_sources.md and sources/book1/BRIEFS.md (g9/01).
\chapter{Web Pages: HTML and CSS}\label{ch:g9:html-and-css}

Lina saves the library's \omterm{def:g5:networks-first-look:page}{web page} on her laptop and opens the \omterm{def:g3:files-and-folders:file}{file} in a \omterm{def:g2:using-a-computer:word-processor}{text
editor} instead of the \omterm{def:g5:networks-first-look:page}{browser}. There is no picture of the page: only words,
and many small words between angle brackets, such as \code{<h1>} and
\code{</p>}. Those small words tell the \omterm{def:g5:networks-first-look:page}{browser} what each part of the page
is. This chapter reads them.

\begin{recall}
A \omterm{def:g5:networks-first-look:page}{web page} is a \omterm{def:g3:files-and-folders:file}{file} kept by a \omterm{def:g5:networks-first-look:server}{server}; a \omterm{def:g5:networks-first-look:page}{browser} asks for it and shows it; a
\omterm{def:g5:networks-first-look:web}{link} holds the address of another page (\cref{ch:g5:networks-first-look}). A
\omterm{def:g3:files-and-folders:file}{file} has a name and an \omterm{def:g3:files-and-folders:extension}{extension} that tells what it holds
(\cref{ch:g3:files-and-folders}). A colour can be written \texttt{\#RRGGBB},
three \omterm{def:g4:memory-and-storage:byte}{bytes} in \omterm{def:g8:binary-and-hexadecimal:hexadecimal}{hexadecimal} (\cref{ch:g8:binary-and-hexadecimal}).
\end{recall}

\section{A web page is a text file}

\begin{definition}[HTML]\label{def:g9:html-and-css:html}
\emph{HTML}\index{HTML} (HyperText Markup Language) is the language in which
\omterm{def:g5:networks-first-look:page}{web pages} are written. An HTML \omterm{def:g3:files-and-folders:file}{file} is a \omterm{def:g8:files-and-data-tables:text-file}{text file}, with the \omterm{def:g3:files-and-folders:extension}{extension}
\code{.html}, in which marks around the text say what each part is: a
heading, a paragraph, a \omterm{def:g7:lists:list}{list}, a \omterm{def:g5:networks-first-look:web}{link}.
\end{definition}

Here is the whole page Lina opened:

\omcode{code/grade-9/01-html-and-css/listings/first_page.html}{1}{18}{The library's
page, as text.}\label{lst:g9:html-and-css:first}

\begin{omfigure}
\begin{tikzpicture}[font=\small]
  \draw[black!60, fill=black!8, rounded corners=3pt] (0,0) rectangle (9,0.6);
  \foreach \x/\c in {0.3/ocRed, 0.6/ocYellow, 0.9/ocMint} { \fill[\c] (\x,0.3) circle (0.09); }
  \draw[black!50, fill=white, rounded corners=2pt] (1.4,0.12) rectangle (8.7,0.48);
  \node[anchor=west, font=\scriptsize\ttfamily, black!70] at (1.5,0.3) {first\_page.html};
  \draw[black!60, fill=white] (0,-6.2) rectangle (9,0);
  \node[anchor=west, font=\Large\bfseries, ocMint!60!black] at (0.3,-0.6) {Brenmoor Library};
  \node[anchor=west] at (0.3,-1.3) {Open every day \emph{except} Monday.};
  \node[anchor=west, font=\large\bfseries] at (0.3,-2.0) {This week};
  \node[anchor=west] at (0.5,-2.6) {\textbullet\ Story time on Wednesday};
  \node[anchor=west] at (0.5,-3.1) {\textbullet\ \textbf{Chess club on Saturday}};
  \node[anchor=west] at (0.3,-3.8) {See the \textcolor{ocBlue}{\underline{list of books}}.};
  \draw[black!50, fill=black!5] (0.3,-6.0) rectangle (3.3,-4.3);
  \node[align=center, font=\scriptsize, black!60] at (1.8,-5.15) {the photo\\ \code{library.jpg}};
\end{tikzpicture}

{\small The same \omterm{def:g3:files-and-folders:file}{file} in a \omterm{def:g5:networks-first-look:page}{browser} window, drawn from the \omterm{def:g9:html-and-css:html}{HTML} and its style
sheet: the marks have disappeared, and the text is laid out as headings, a
\omterm{def:g7:lists:list}{list}, a \omterm{def:g5:networks-first-look:web}{link} and an image.}
\end{omfigure}

\section{Elements and tags}

\begin{definition}[Element, tag]\label{def:g9:html-and-css:element}\label{def:g9:html-and-css:tag}
An \emph{element}\index{element (of HTML)} is one part of a page. It is
written as a \emph{start tag}\index{tag (HTML)} such as \code{<p>}, then its
content, then an \emph{end tag} such as \code{</p>}, with a slash. A few
elements, such as an image, have no content and no end tag.
\end{definition}

\begin{definition}[Attribute]\label{def:g9:html-and-css:attribute}
An \emph{attribute}\index{attribute (of HTML)} is a name and a \omterm{def:g6:variables-and-input:value}{value} written
inside a \omterm{def:g9:html-and-css:element}{start tag}, such as \code{href="books.html"}; it gives more
information about the \omterm{def:g9:html-and-css:element}{element}.
\end{definition}

% ledger: g9:html-void
The \omterm{def:g9:html-and-css:element}{elements} of the library page are among the most used:

\begin{center}\small
\begin{tabular}{@{}ll@{}}
\hline
\omterm{def:g9:html-and-css:element}{Element} & What it is \\
\hline
\code{h1}, \code{h2} & a heading, then a smaller heading (down to \code{h6}) \\
\code{p} & a paragraph \\
\code{em} & words to stress, usually shown in italics \\
\code{ul}, \code{li} & a \omterm{def:g7:lists:list}{list} with bullets, and one \omterm{def:g7:lists:list}{item} of it \\
\code{a} & a \omterm{def:g5:networks-first-look:web}{link}; its \omterm{def:g9:html-and-css:attribute}{attribute} \code{href} gives the address \\
\code{img} & an image; \code{src} names the \omterm{def:g3:files-and-folders:file}{file}, \code{alt} describes it \\
\code{table}, \code{tr}, \code{th}, \code{td} & a table, a row, a header cell, a cell \\
\code{html}, \code{head}, \code{body} & the whole page, its information, what is shown \\
\hline
\end{tabular}
\end{center}

The \code{head} holds information about the page: its \code{title}, shown on
the \omterm{def:g5:networks-first-look:page}{browser}'s tab, and a \code{link} to its style sheet. The \code{body} holds
everything shown in the window. The first line, \code{<!DOCTYPE html>}, only
tells the \omterm{def:g5:networks-first-look:page}{browser} that the \omterm{def:g3:files-and-folders:file}{file} is \omterm{def:g9:html-and-css:html}{HTML}. A table is written row by row:

\omcode{code/grade-9/01-html-and-css/listings/table.html}{1}{14}{A small table of
the town's places.}\label{lst:g9:html-and-css:table}

The \code{alt} text of an image is read aloud to people who cannot see the
screen, and shown when the image does not load: a page should always have one.

\section{The page as a tree}

\begin{definition}[Document tree]\label{def:g9:html-and-css:tree}
The \omterm{def:g9:html-and-css:element}{elements} of a page sit inside one another: \code{html} holds \code{head}
and \code{body}, \code{body} holds the paragraphs, a paragraph holds a \omterm{def:g5:networks-first-look:web}{link}.
Drawn with each \omterm{def:g9:html-and-css:element}{element} below the one that holds it, they make the
\emph{document tree}\index{document tree} of the page.
\end{definition}

A \omterm{def:g4:programs-and-operating-system:program}{program} can find the tree. It looks for every \code{<}, cuts out the tag up
to the next \code{>}, and keeps a count of the open \omterm{def:g9:html-and-css:element}{elements}, the
\code{depth}: a \omterm{def:g9:html-and-css:element}{start tag} opens one more, an end tag closes one. Each \omterm{def:g9:html-and-css:element}{element}
name is printed with two spaces per open \omterm{def:g9:html-and-css:element}{element} in front of it.

\omcode{code/grade-9/01-html-and-css/listings/html_tree.py}{1}{33}{Tags, \omterm{def:g9:html-and-css:element}{element}
names and the tree.}\label{lst:g9:html-and-css:html-tree}

\Needspace*{9\baselineskip}
\omcode{code/grade-9/01-html-and-css/listings/show_tree.py}{1}{6}{The tree of the
library page.}\label{lst:g9:html-and-css:show-tree}

\omoutput{code/grade-9/01-html-and-css/outputs/show_tree.txt}{Indentation shows
what is inside what; at the end, no \omterm{def:g9:html-and-css:element}{element} is left
open.}\label{out:g9:html-and-css:show-tree}

\code{find} returns the position of a text inside another, or $-1$ when it is
not there; \code{find("<", end)} starts looking at position \code{end}.
\code{continue} skips the rest of the \omterm{def:g4:algorithms-repeat-and-decide:loop}{loop}'s body and goes on with the next
tag. \code{link} and \code{img} are void \omterm{def:g9:html-and-css:element}{elements}, with no end tag, so they do
not open anything.

\begin{omfigure}
\begin{forest}
  for tree={font=\ttfamily\small, draw=black!50, rounded corners=2pt, fill=ocBlue!8,
    inner sep=2pt, l sep=4mm, s sep=2mm, edge={black!50}}
  [html
    [head [title] [link]]
    [body [h1] [p [em]] [h2] [ul [li] [li]] [p [a]] [img]]
  ]
\end{forest}

{\small The \omterm{def:g9:html-and-css:tree}{document tree} of the library page: the same 15 \omterm{def:g9:html-and-css:element}{elements} as the
indented \omterm{def:g1:what-is-a-computer:input}{output} of the \omterm{def:g4:programs-and-operating-system:program}{program}.}
\end{omfigure}

The tree only comes out right if every end tag closes the \omterm{def:g9:html-and-css:element}{element} opened last.
A missing end tag shifts everything after it:

\Needspace*{11\baselineskip}
\omcode{code/grade-9/01-html-and-css/listings/unclosed.py}{1}{8}{A page where
\code{em} is never closed.}\label{lst:g9:html-and-css:unclosed}

\Needspace*{10\baselineskip}
\omoutput{code/grade-9/01-html-and-css/outputs/unclosed.txt}{The heading and the
second paragraph end up inside the first paragraph.}\label{out:g9:html-and-css:unclosed}

The \code{</p>} closes \code{em} instead of \code{p}, so the paragraph stays
open, and the heading lands inside it. \omterm{def:g5:networks-first-look:page}{Browsers} try hard to guess what was
meant, but a page whose tags pair up properly leaves nothing to guess. The
\omterm{def:g4:programs-and-operating-system:program}{program} can also count the \omterm{def:g9:html-and-css:element}{elements} of each kind:

\omcode{code/grade-9/01-html-and-css/listings/count_tags.py}{1}{13}{Counting the
\omterm{def:g9:html-and-css:element}{elements}.}\label{lst:g9:html-and-css:count}

\Needspace*{12\baselineskip}
\omoutput[1-5,13-15]{code/grade-9/01-html-and-css/outputs/count_tags.txt}{15
\omterm{def:g9:html-and-css:element}{elements}.}\label{out:g9:html-and-css:count}

\begin{history}[The first web page]
% ledger: g5:web-first-server
% ledger: g9:web-first-page
The first web \omterm{def:g5:networks-first-look:server}{server} ran at CERN at the end of 1990
(\cref{ch:g5:networks-first-look}). Its address was \code{info.cern.ch}, and
its first page described the \omterm{def:g5:networks-first-look:web}{World Wide Web} project itself, with \omterm{def:g5:networks-first-look:web}{links} to
pages on how to set up a \omterm{def:g5:networks-first-look:server}{server} of one's own.
\end{history}

\section{Content apart from appearance}

\begin{definition}[CSS, CSS rule, selector]\label{def:g9:html-and-css:css}\label{def:g9:html-and-css:rule}\label{def:g9:html-and-css:selector}
\emph{CSS}\index{CSS} (Cascading Style Sheets) is the language that says how
the \omterm{def:g9:html-and-css:element}{elements} of a page look. A \emph{CSS rule}\index{CSS rule} is a
\emph{selector}\index{selector (CSS)}, which says which \omterm{def:g9:html-and-css:element}{elements} it applies
to, followed by properties and \omterm{def:g6:variables-and-input:value}{values} between curly brackets.
\end{definition}

\begin{definition}[Class of an element]\label{def:g9:html-and-css:class}
The \omterm{def:g9:html-and-css:attribute}{attribute} \code{class} gives an \omterm{def:g9:html-and-css:element}{element} a name chosen by the author, its
\emph{class}\index{class (of an HTML element)}; the \omterm{def:g9:html-and-css:css}{selector}
\code{li.new} applies only to the \code{li} \omterm{def:g9:html-and-css:element}{elements} of class \code{new}.
\end{definition}

The \omterm{def:g9:html-and-css:html}{HTML} says what each part is; a separate \omterm{def:g9:html-and-css:css}{CSS} \omterm{def:g3:files-and-folders:file}{file}, the style sheet, says
how it looks. The \code{link} in the \code{head} of the library page names
\code{style.css}:

\omcode{code/grade-9/01-html-and-css/listings/style.css}{1}{20}{The style sheet of
the library page.}\label{lst:g9:html-and-css:style}

The \omterm{def:g9:html-and-css:css}{selector} \code{body} applies to the body and, unless a rule says
otherwise, to everything inside it: hence the font of the whole page. The
colours are the \omterm{def:g8:binary-and-hexadecimal:hexadecimal}{hexadecimal} colours of \cref{ch:g8:binary-and-hexadecimal};
sizes are in \code{px}, \omterm{def:g5:text-pictures-sound-as-numbers:pixel}{pixels} of the screen.

\section{Boxes}

\begin{definition}[Box model]\label{def:g9:html-and-css:box-model}
In the \emph{box model}\index{box model}, the \omterm{def:g5:networks-first-look:page}{browser} draws every \omterm{def:g9:html-and-css:element}{element} as
a rectangle: its content, surrounded by an inner space, the \emph{padding};
then a \emph{border}; then an outer space, the \emph{margin}, that keeps
other boxes away.
\end{definition}

% ledger: g9:css-box
In \code{style.css}, a paragraph has content 300 \omterm{def:g5:text-pictures-sound-as-numbers:pixel}{pixels} wide, 10 \omterm{def:g5:text-pictures-sound-as-numbers:pixel}{pixels} of
padding and a border of 2 \omterm{def:g5:text-pictures-sound-as-numbers:pixel}{pixels} on each side: the box is
$300 + 2 \times 10 + 2 \times 2 = 324$ \omterm{def:g5:text-pictures-sound-as-numbers:pixel}{pixels} wide, and its margin of 20
\omterm{def:g5:text-pictures-sound-as-numbers:pixel}{pixels} keeps the next box at a distance.

\begin{omfigure}
\begin{tikzpicture}[x=0.022cm, y=0.022cm, font=\scriptsize]
  \fill[ocYellow!25] (0,0) rectangle (364,140);
  \fill[ocBlue!70] (20,20) rectangle (344,120);
  \fill[ocMint!25] (22,22) rectangle (342,118);
  \fill[white, draw=black!40, dashed] (32,32) rectangle (332,108);
  \node at (182,70) {content: 300 px wide};
  \node[anchor=north] at (182,30) {padding 10};
  \node[anchor=south] at (182,121) {border 2};
  \node[anchor=south] at (182,141) {margin 20};
  \draw[{Stealth[length=1.5mm]}-{Stealth[length=1.5mm]}] (20,-8) -- node[below] {324 px} (344,-8);
  \draw[{Stealth[length=1.5mm]}-{Stealth[length=1.5mm]}] (0,-30) -- node[below] {364 px with the margins} (364,-30);
\end{tikzpicture}

{\small The box of a paragraph with \code{style.css}: the content, the
padding, the border (blue) and the margin, to scale.}
\end{omfigure}

Changing the style sheet changes the look of the page without touching the
\omterm{def:g9:html-and-css:html}{HTML}. Here is a second one, dark, with no border around paragraphs:

\omcode{code/grade-9/01-html-and-css/listings/two_styles.css}{1}{21}{A second style
sheet for the same page.}\label{lst:g9:html-and-css:two-styles}

\begin{omfigure}
\begin{tikzpicture}[font=\scriptsize]
  \begin{scope}
    \draw[black!60, fill=white] (0,0) rectangle (6.0,-4.2);
    \node[anchor=west, font=\scriptsize\bfseries\sffamily, ocMint!60!black] at (0.2,-0.4) {Brenmoor Library};
    \draw[ocBlue, thick] (0.25,-0.8) rectangle (5.75,-1.4);
    \node[anchor=west, font=\scriptsize\sffamily] at (0.4,-1.1) {Open every day \emph{except} Monday.};
    \node[anchor=west, font=\scriptsize\bfseries\sffamily] at (0.2,-1.8) {This week};
    \node[anchor=west, font=\scriptsize\sffamily] at (0.3,-2.25) {\textbullet\ Story time on Wednesday};
    \node[anchor=west, font=\scriptsize\bfseries\sffamily] at (0.3,-2.65) {\textbullet\ Chess club on Saturday};
    \node[anchor=north] at (3.0,-4.3) {\code{style.css}};
  \end{scope}
  \begin{scope}[xshift=6.4cm]
    \definecolor{dk}{HTML}{222222}\definecolor{orng}{HTML}{FF8800}
    \draw[black!60, fill=dk] (0,0) rectangle (6.0,-4.2);
    \node[font=\bfseries, orng] at (3.0,-0.4) {Brenmoor Library};
    \node[anchor=west, white] at (0.2,-1.1) {Open every day \emph{except} Monday.};
    \node[anchor=west, font=\bfseries, white] at (0.2,-1.8) {This week};
    \node[anchor=west, white] at (0.3,-2.25) {\textbullet\ Story time on Wednesday};
    \node[anchor=west, orng] at (0.3,-2.65) {\textbullet\ Chess club on Saturday};
    \node[anchor=north] at (3.0,-4.3) {\code{two_styles.css}};
  \end{scope}
\end{tikzpicture}

{\small The top of the same \omterm{def:g9:html-and-css:html}{HTML} page with the two style sheets: only the \omterm{def:g9:html-and-css:css}{CSS}
changed.}
\end{omfigure}

\section{Exercises}

\begin{exercise}[$\star$]\label{exo:g9:html-and-css:1}
What does \code{show_tree} print for this \omterm{def:g9:html-and-css:html}{HTML}?

\smallskip\centerline{\code{<ul><li><a href="x.html">x</a></li><li>y</li></ul>}}
\end{exercise}

\begin{exercise}[$\star$]\label{exo:g9:html-and-css:2}
In the library page, which \omterm{def:g9:html-and-css:element}{elements} does each rule of \code{style.css}
apply to?
\end{exercise}

\begin{exercise}[$\star$]\label{exo:g9:html-and-css:3}
With \code{two_styles.css}, how wide is the box of a paragraph, border
included? And with its margins?
\end{exercise}

\begin{exercise}[$\star$]\label{exo:g9:html-and-css:4}
Write a page with a title, a heading, a \omterm{def:g7:lists:list}{list} of two of your favourite places
and a \omterm{def:g5:networks-first-look:web}{link} to \code{https://www.example.org}.
\end{exercise}

\begin{exercise}[$\star\star$]\label{exo:g9:html-and-css:5}
\code{<p>An <em>important</p></em> day} closes its tags in the wrong order.
Draw the tree it should have, and fix it.
\end{exercise}

\begin{exercise}[$\star\star$]\label{exo:g9:html-and-css:6}
A page shows a map with \code{<img src="map.png">}. What is missing, who
needs it, and how do you fix it?
\end{exercise}

\begin{exercise}[$\star\star$]\label{exo:g9:html-and-css:7}
Write an \omterm{def:g9:html-and-css:html}{HTML} table with two columns, the name and the kind of four of the
town's places: Library, Market Hall, River Park and one more of your choice.
\end{exercise}

\begin{exercise}[$\star\star$]\label{exo:g9:html-and-css:8}
Write a style sheet that colours every \code{h2} heading in royal blue,
\texttt{\#4169E1}.
\end{exercise}

\begin{exercise}[$\star\star$]\label{exo:g9:html-and-css:9}
How many \omterm{def:g9:html-and-css:element}{elements} does \code{table.html} have, and what is the depth of its
tree (the number of levels, the top one counted as 1)?
\end{exercise}

\begin{exercise}[$\star\star\star$]\label{exo:g9:html-and-css:10}
Write a \omterm{def:g7:functions:function}{function} \code{links(text)} that returns the address of every \omterm{def:g5:networks-first-look:web}{link} of
a page, using \code{tags} and \code{name_of} from
\cref{lst:g9:html-and-css:html-tree}.
\end{exercise}

\begin{exercise}[$\star\star\star$]\label{exo:g9:html-and-css:11}
Show that in a page whose tags pair up properly, every end tag closes the
\omterm{def:g9:html-and-css:element}{element} opened most recently and still open, as closing brackets do in
$(\,[\,]\,(\,)\,)$. Give a page where this fails, and say what
\code{show_tree} does with it.
\end{exercise}

\begin{exercise}[$\star\star\star$]\label{exo:g9:html-and-css:12}
Write a \omterm{def:g7:functions:function}{function} \code{html_table(filename)} that reads a \omterm{def:g8:files-and-data-tables:table}{CSV file}
(\cref{ch:g8:files-and-data-tables}) and returns the text of an \omterm{def:g9:html-and-css:html}{HTML} table,
with the \omterm{def:g8:files-and-data-tables:table}{header row} in \code{th} cells.
\end{exercise}

\section{Problem: The Tree of a Page}

\begin{problem}[{Weekend problem --- the tree of a page}]\label{pb:g9:html-and-css:1}
The weather station of Hilltop publishes a small page:

\omcode{code/grade-9/01-html-and-css/listings/weather_page.html}{1}{19}{The weather
page.}\label{lst:g9:html-and-css:weather}

\medskip\noindent\textbf{Part I --- By hand.}
\begin{enumerate}
  \item Draw the tree of the page.
  \item How many \omterm{def:g9:html-and-css:element}{elements} does it have?
  \item What is the depth of its tree, the top level counted as 1? Which
    \omterm{def:g9:html-and-css:element}{elements} are the deepest?
\end{enumerate}

\medskip\noindent\textbf{Part II --- By program.}

\omcode{code/grade-9/01-html-and-css/listings/page_stats.py}{1}{38}{The number of
\omterm{def:g9:html-and-css:element}{elements} and the depth of a page.}\label{lst:g9:html-and-css:stats}

\begin{enumerate}[resume]
  \item Why is \code{deepest} updated with \code{depth + 1} before
    \code{depth} itself changes?
  \item Why does an \code{img} count as an \omterm{def:g9:html-and-css:element}{element}, but never change
    \code{depth}?
  \item Run the \omterm{def:g4:programs-and-operating-system:program}{program}. Do its numbers for the weather page agree with
    Part I?
  \item What does \code{list_page(3)} return? Write its tree.
\end{enumerate}

\medskip\noindent\textbf{Part III --- A growing list.}
\begin{enumerate}[resume]
  \item How do the count and the depth change from a \omterm{def:g7:lists:list}{list} of 1 \omterm{def:g7:lists:list}{item} to a \omterm{def:g7:lists:list}{list}
    of 10?
  \item Write the count as a formula in $n$, the number of \omterm{def:g7:lists:list}{items}.
  \item Why does the depth not change?
  \item State the result.
\end{enumerate}
\end{problem}
